\documentclass[12pt,a4paper]{article}

% --- Pacotes básicos ---
\usepackage[utf8]{inputenc}
\usepackage[T1]{fontenc}
\usepackage[brazil]{babel}
\usepackage{geometry}
\usepackage{setspace}
\usepackage{graphicx}
\usepackage{hyperref}

\geometry{margin=2.5cm}
\onehalfspacing

\title{\textbf{Relatório Técnico}\\[0.4em]
\large Visualizador Interativo de Modelos 3D (OBJ/MTL) utilizando WebGL}
\author{Disciplina de Computação Gráfica}
\date{\today}

\begin{document}

\maketitle

\begin{center}
\textbf{Integrantes do Grupo}

Henrique Froeder \\
Nicolas Castro \\
Jean Kronbauer \\
Igor Paslauski Pedroso de Oliveira
\end{center}

\begin{abstract}
Este trabalho apresenta o desenvolvimento de um visualizador de modelos 3D executado diretamente no navegador. A aplicação foi construída utilizando HTML, CSS e JavaScript puro, com renderização via WebGL.

O sistema permite carregar arquivos no formato OBJ e MTL, aplicar transformações no modelo em tempo real e visualizar a cena em diferentes modos, como sólido e wireframe. Também são exibidas informações sobre a malha 3D, como número de vértices, arestas e faces.
\end{abstract}

\tableofcontents
\newpage

%===========================================================
\section{Introdução}
%===========================================================

A computação gráfica permite a criação e visualização de objetos tridimensionais em ambientes digitais. Um dos formatos mais utilizados para representar modelos 3D é o formato OBJ, que descreve a geometria do objeto. Já o formato MTL define suas propriedades visuais, como cores.

Neste projeto, foi desenvolvido um visualizador 3D simples, com o objetivo de demonstrar conceitos fundamentais como transformações geométricas, projeção e renderização em tempo real.

%===========================================================
\section{Objetivos}
%===========================================================

\subsection{Objetivo Geral}
Desenvolver um visualizador 3D funcional no navegador.

\subsection{Objetivos Específicos}
\begin{itemize}
    \item Ler arquivos OBJ e MTL
    \item Converter modelos para triângulos renderizáveis
    \item Aplicar transformações (rotação, escala e translação)
    \item Implementar projeções perspectiva e ortográfica
    \item Exibir modelos em modo sólido e wireframe
    \item Mostrar estatísticas da malha (vértices, arestas e faces)
\end{itemize}

%===========================================================
\section{Arquitetura do Sistema}
%===========================================================

O sistema foi dividido em módulos para facilitar organização e manutenção:

\begin{itemize}
    \item \textbf{index.html}: estrutura da aplicação
    \item \textbf{style.css}: estilização da interface
    \item \textbf{main.js}: controle geral e interação do usuário
    \item \textbf{math3d.js}: cálculos matemáticos (matrizes e vetores)
    \item \textbf{objParser.js}: leitura de arquivos OBJ
    \item \textbf{mtlParser.js}: leitura de arquivos MTL
    \item \textbf{renderer.js}: renderização utilizando WebGL
\end{itemize}

Essa divisão permite que cada parte do sistema tenha uma responsabilidade específica.

%===========================================================
\section{Leitura dos Arquivos}
%===========================================================

\subsection{Arquivo OBJ}

O arquivo OBJ contém a geometria do modelo 3D. Nele são definidos:

\begin{itemize}
    \item Vértices (v)
    \item Normais (vn)
    \item Faces (f)
\end{itemize}

As faces são convertidas em triângulos para facilitar a renderização.

\subsection{Arquivo MTL}

O arquivo MTL define os materiais do objeto, como cores. Neste projeto, foi utilizado apenas o parâmetro de cor difusa (Kd).

%===========================================================
\section{Renderização}
%===========================================================

A renderização é feita utilizando WebGL, que permite desenhar gráficos 3D diretamente no navegador.

\subsection{Iluminação}

Foi implementada iluminação difusa simples (modelo de Lambert), que considera a incidência de luz sobre a superfície do objeto.

\subsection{Projeções}

O sistema suporta dois tipos de projeção:

\begin{itemize}
    \item Perspectiva (mais realista)
    \item Ortográfica (sem profundidade visual)
\end{itemize}

%===========================================================
\section{Transformações}
%===========================================================

O usuário pode interagir com o modelo aplicando:

\begin{itemize}
    \item Rotação
    \item Escala
    \item Translação
\end{itemize}

Essas transformações são realizadas utilizando matrizes 4x4.

%===========================================================
\section{Interface}
%===========================================================

A interface foi desenvolvida com foco em simplicidade e usabilidade. Ela permite:

\begin{itemize}
    \item Carregar arquivos OBJ e MTL
    \item Alternar entre modo sólido e wireframe
    \item Visualizar informações do modelo
\end{itemize}

%===========================================================
\section{Limitações}
%===========================================================

\begin{itemize}
    \item Não há suporte para texturas
    \item Iluminação simplificada
    \item Alguns modelos complexos podem não ser renderizados corretamente
\end{itemize}

%===========================================================
\section{Conclusão}
%===========================================================

O projeto atingiu seu objetivo de implementar um visualizador 3D funcional e didático. Ele permite compreender, na prática, conceitos importantes da computação gráfica, como leitura de arquivos, transformações e renderização.

Além disso, a estrutura modular facilita futuras melhorias, como suporte a texturas e iluminação mais avançada.

%===========================================================
\section{Referências}
%===========================================================

Shreiner, D. et al. \textit{OpenGL Programming Guide}. Addison-Wesley.

\end{document}